% Generated draft: requires mathematical and source review.
\documentclass[11pt]{article}
\usepackage[margin=25mm]{geometry}
\usepackage{fontspec}
\setmainfont{Latin Modern Roman}
\usepackage{amsmath,amssymb,mathtools,unicode-math}
\setmathfont{Latin Modern Math}
\usepackage{enumitem,needspace,xurl,hyperref}
\setlength{\parindent}{0pt}
\setlength{\parskip}{5pt}
\setlength{\emergencystretch}{3em}
\newcommand{\N}{\mathbb{N}}
\newcommand{\Z}{\mathbb{Z}}
\newcommand{\Q}{\mathbb{Q}}
\newcommand{\R}{\mathbb{R}}
\newcommand{\C}{\mathbb{C}}
\newcommand{\F}{\mathbb{F}}
\newcommand{\A}{\mathbb{A}}
\newcommand{\PP}{\mathbb{P}}
\newcommand{\E}{\mathbb{E}}
\newcommand{\eps}{\varepsilon}
\newcommand{\dd}{\,\mathrm{d}}
\begin{document}
\section*{UnsolvedMath ID: 30000824}
\textbf{Source problem code:} catalogue entry
\subsection{Definitions and mathematical statement}
For context over $\C$, let $V_{N,d}=H^0(\PP^N,\mathcal O(d))$ be the vector space of degree-$d$ homogeneous forms, and let $\mathcal R_{N,d}\subset\PP(V_{N,d})$ be the locus of completely reducible forms, meaning products of $d$ linear forms under the convention allowing repeated factors. A pencil is a projective line in this parameter space; more generally a linear system is a projective linear subspace $L\subset\PP(V_{N,d})$. Completely reducible members correspond to $L\cap\mathcal R_{N,d}$.

The source reports that a pencil of plane curves has at most four completely reducible fibers, in its intended setting, and asks for two extensions:
\begin{enumerate}
\item analogous upper bounds for linear systems in higher dimensions;
\item analogous upper bounds when the parameter locus is a curve $C\subset\PP(V_{N,d})$ of degree greater than one, rather than a pencil line.
\end{enumerate}
For the latter, a natural counting target is the number of distinct points of $C\cap\mathcal R_{N,d}$ when this intersection is finite. The degree of the parameter curve is distinct from the degree $d$ of the represented hypersurfaces. For the former, ``higher dimensions'' could refer to the ambient projective space, the parameter dimension, or both; a positive-dimensional reducible-member locus requires a different type of bound from a count of fibers.

The four-fiber assertion cannot be read without nondegeneracy hypotheses. For example, the plane pencil spanned by $x^d$ and $y^d$ consists entirely of products of linear forms over $\C$, so arbitrary pencils need not have finitely many completely reducible members. Likewise a parameter curve contained in $\mathcal R_{N,d}$ has infinitely many. The extract does not provide the hypotheses excluding such pencils or curves, fixed components or composite systems, nor the permitted multiplicities, field, dependence of bounds or intersection-count convention. These are source gaps; no bound for all unrestricted systems or a guessed constant is asserted.
\subsection{Short English statement}
Extend the source's four-completely-reducible-fiber bound to higher-dimensional systems and to nonlinear parameter curves of hypersurfaces. The intended nondegeneracy and counting hypotheses are missing; unrestricted families can consist entirely of completely reducible forms.
\subsection{Sources}
\begin{itemize}
\item Supplied report: \emph{Mini-Workshop: Topology of closed one-forms and Cohomology Jumping Loci} (2008), Oberwolfach Reports, DOI 10.4171/owr/2007/40, pp. 2321--2360; proposed year 2007. \url{https://doi.org/10.4171/owr/2007/40}.
\item Local UnsolvedMath registry, permanent ID 30000824, catalogue entry; created 25 August 2026, updated 16 September 2026. It labels the statement exact and reports comparison with the official report. The supplied 21 August 2026 literature triage labels the broad upper-bound direction open and reports partial arrangement results rather than a general answer. No additional external reference is supplied, and no independent status assessment is asserted here.
\end{itemize}
\end{document}
